\section{The shift spectrum of a smooth determinant packet}
\label{sec:tpc46-shift-spectrum}

We study the completed Hilbert packet
\eqref{eq:tpc46-Hilbert-packet}.  This is a Fourier model for the
equation \(pr-mj=h\); it is not yet an assertion that the actual TPC
packet has been constructed uniformly for every \(h\).  Let \(J\ge2\)
and suppose that its source support satisfies

\[
 0<M_-\le m\le M_+,
\]

The coefficients and vectors have finite arithmetic support and
\(w\in\mathcal S(\R)\).  The vectors in
\eqref{eq:tpc46-Hilbert-packet} may depend jointly and arbitrarily on
\((p,r,m)\), but they do not depend on \(j\).  In particular, the
residual-label choice
\begin{equation}
 v_{p,r,m}=e_{d(m)r}
 \quad\text{in }\ell^2(\N)
 \label{eq:tpc46-residual-basis-example}
\end{equation}
is permitted.  The sequence \((F_h)_h\) belongs to both
\(\ell^1(\Z;\Hh)\) and \(\ell^2(\Z;\Hh)\).  We use the shift Fourier
transform
\begin{equation}
 \widehat F(\alpha):=
 \sum_{h\in\Z}F_h\e(-h\alpha),
 \qquad \alpha\in\T.
 \label{eq:tpc46-shift-Fourier-transform}
\end{equation}

\begin{theorem}[Hilbert-valued orbit kernel]
\label{thm:tpc46-Hilbert-orbit-kernel}
For every \(\alpha\in\T\),
\begin{align}
 \widehat F(\alpha)
 &=\sum_{p,r,m}
   z_{p,r,m}v_{p,r,m}\e(-pr\alpha)K_m(\alpha),
 \label{eq:tpc46-shift-kernel-factorization}\\
 K_m(\alpha)
 &:=\sum_{j\in\Z}w(j/J)\e(mj\alpha)
   =J\sum_{k\in\Z}\widehat w\bigl(J(k-m\alpha)\bigr).
 \label{eq:tpc46-Poisson-orbit-kernel}
\end{align}
Moreover,
\begin{equation}
 \sum_{h\in\Z}\|F_h\|_{\Hh}^2
 =\int_{\T}\|\widehat F(\alpha)\|_{\Hh}^2\,d\alpha.
 \label{eq:tpc46-Hilbert-shift-Parseval}
\end{equation}
\end{theorem}

\begin{proof}
Absolute convergence allows the order of summation to be changed.
Substituting \(h=pr-mj\) into
\eqref{eq:tpc46-shift-Fourier-transform} gives the first identity.
Poisson summation applied to \(x\mapsto w(x/J)\) gives
\[
 \sum_{j\in\Z}w(j/J)\e(mj\alpha)
 =\sum_{k\in\Z}J\widehat w\bigl(J(k-m\alpha)\bigr),
\]
which proves the second.  The last assertion is the standard
Hilbert-valued Parseval identity; it also follows by applying scalar
Parseval after pairing with an orthonormal basis of the finite
dimensional span of the vectors \(v_{p,r,m}\).
\end{proof}

The factor \(K_m\) contains all shift frequencies contributed by the
smooth orbit variable.  Neither the primality of \(p\) nor any
arithmetic cancellation in the three coefficient slots is used in
the following theorem.

\begin{theorem}[Exact mesoscopic shift-spectral gap]
\label{thm:tpc46-exact-shift-spectral-gap}
Assume that, for some \(\sigma>0\),
\begin{equation}
 \operatorname{supp}\widehat w\subset[-\sigma,\sigma]
 \qquad\text{and}\qquad J>\sigma.
 \label{eq:tpc46-bandlimited-orbit-profile}
\end{equation}
Then
\begin{equation}
 \widehat F(\alpha)=0
 \label{eq:tpc46-exact-gap-vanishing}
\end{equation}
whenever
\begin{equation}
 \boxed{
 \frac{\sigma}{M_-J}
 <\|\alpha\|_{\T}<
 \frac{1-\sigma/J}{M_+}.}
 \label{eq:tpc46-exact-gap-annulus}
\end{equation}
The assertion is vacuous if the displayed interval is empty.
\end{theorem}

\begin{proof}
Formula \eqref{eq:tpc46-Poisson-orbit-kernel} and
\eqref{eq:tpc46-bandlimited-orbit-profile} show that
\begin{equation}
 K_m(\alpha)\ne0
 \quad\Longrightarrow\quad
 \|m\alpha\|_{\T}\le\frac{\sigma}{J}.
 \label{eq:tpc46-kernel-support-condition}
\end{equation}
Put \(t=\|\alpha\|_{\T}\) and choose a representative with
\(|\alpha|=t\).  Under \eqref{eq:tpc46-exact-gap-annulus}, every
\(m\in[M_-,M_+]\) satisfies
\[
 \frac{\sigma}{J}<mt<1-\frac{\sigma}{J}.
\]
Thus the distance of \(m\alpha\) from both \(0\) and the nearest
nonzero integer is greater than \(\sigma/J\).  This contradicts
\eqref{eq:tpc46-kernel-support-condition}, so every \(K_m(\alpha)\)
in \eqref{eq:tpc46-shift-kernel-factorization} vanishes.
\end{proof}

\begin{corollary}[Fixed-progression shift gap]
\label{cor:tpc46-progression-shift-gap}
Let \(H_0\in\N\), let \(a\) be a residue class modulo \(H_0\), and
replace \eqref{eq:tpc46-Hilbert-packet} by
\begin{equation}
 F_h^{a,H_0}
 :=\sum_{p,r,m}\sum_{\substack{j\in\Z\\j\equiv a\!\!\pmod{H_0}}}
 z_{p,r,m}w(j/J)v_{p,r,m}\one_{pr-mj=h}.
 \label{eq:tpc46-progression-Hilbert-packet}
\end{equation}
Assume
\(\operatorname{supp}\widehat w\subset[-\sigma,\sigma]\) and
\(J>H_0\sigma\).  Then
\begin{equation}
 \widehat F^{a,H_0}(\alpha)=0
 \label{eq:tpc46-progression-gap-vanishing}
\end{equation}
whenever
\begin{equation}
 \boxed{
 \frac{\sigma}{M_-J}
 <\|\alpha\|_{\T}<
 \frac{H_0^{-1}-\sigma/J}{M_+}.}
 \label{eq:tpc46-progression-gap-annulus}
\end{equation}
For fixed \(H_0\), the endpoint exponent range remains
\(X^{-1+o(1)}<\|\alpha\|_{\T}<X^{-267/400+o(1)}\).
\end{corollary}

\begin{proof}
The progression-restricted orbit kernel is
\begin{align}
 K_{m,a}^{(H_0)}(\alpha)
 &:=
 \sum_{\substack{j\in\Z\\j\equiv a\!\!\pmod{H_0}}}
 w(j/J)\e(mj\alpha)\notag\\
 &=\frac{J}{H_0}\sum_{k\in\Z}
 \e(ak/H_0)
 \widehat w\!\left(
 \frac{J}{H_0}(k-mH_0\alpha)
 \right).
 \label{eq:tpc46-progression-orbit-kernel}
\end{align}
Indeed, write \(j=a+H_0t\) and apply Poisson summation in \(t\).
Thus a nonzero kernel requires
\[
 \|mH_0\alpha\|_{\T}\le\frac{H_0\sigma}{J}.
\]
If \(t=\|\alpha\|_{\T}\) lies in
\eqref{eq:tpc46-progression-gap-annulus}, then every supported \(m\)
satisfies
\[
 \frac{H_0\sigma}{J}<mH_0t<1-\frac{H_0\sigma}{J}.
\]
The same support argument as in
\cref{thm:tpc46-exact-shift-spectral-gap} makes every
\(K_{m,a}^{(H_0)}(\alpha)\) vanish.  The endpoint statement follows
because \(H_0\) is fixed.
\end{proof}

\begin{corollary}[The endpoint spectral no-man's-land]
\label{cor:tpc46-endpoint-shift-gap}
Suppose
\begin{equation}
 M_-=X^{267/400+o(1)},\qquad
 M_+=X^{267/400+o(1)},\qquad
 J=X^{133/400+o(1)}.
 \label{eq:tpc46-endpoint-MJ-scales}
\end{equation}
For a fixed band limit \(\sigma\), the exact gap in
\cref{thm:tpc46-exact-shift-spectral-gap} has the exponent form
\begin{equation}
 X^{-1+o(1)}
 <\|\alpha\|_{\T}<
 X^{-267/400+o(1)}.
 \label{eq:tpc46-endpoint-frequency-gap}
\end{equation}
Thus the zero-frequency orbit envelope, whose width is of order
\((MJ)^{-1}\), is separated from the first lattice aliases, whose
distance from zero is of order \(M^{-1}\).
\end{corollary}

Strict Fourier compactness is convenient for an exact statement, but
is not needed for a power-robust gap.

\begin{theorem}[Rapid decay for a smooth compact orbit profile]
\label{thm:tpc46-rapid-smooth-shift-gap}
Let \(w\in C_c^\infty(\R)\), and suppose that for fixed
\(\mu,\nu>0\) with \(\mu+\nu=1\),
\begin{equation}
 M_-\asymp M_+\asymp X^\mu,
 \qquad J\asymp X^\nu.
 \label{eq:tpc46-general-MJ-scales}
\end{equation}
Fix \(0<\epsilon<\nu\).  For every \(A>0\), uniformly on
\begin{equation}
 X^{-1+\epsilon}
 \le\|\alpha\|_{\T}
 \le X^{-\mu-\epsilon},
 \label{eq:tpc46-shrunken-smooth-gap}
\end{equation}
one has
\begin{align}
 |K_m(\alpha)|
 &\ll_{A,w,\epsilon}JX^{-A\epsilon}
 \qquad(M_-\le m\le M_+),
 \label{eq:tpc46-smooth-kernel-rapid-decay}\\
 \|\widehat F(\alpha)\|_{\Hh}
 &\ll_{A,w,\epsilon}
 JX^{-A\epsilon}
 \sum_{p,r,m}|z_{p,r,m}|\|v_{p,r,m}\|_{\Hh}.
 \label{eq:tpc46-smooth-packet-rapid-decay}
\end{align}
Replacing \(A\) by a larger constant absorbs any fixed polynomial
size of the coefficient support.
\end{theorem}

\begin{proof}
Let \(t=\|\alpha\|_{\T}\).  Uniformly in the range
\eqref{eq:tpc46-shrunken-smooth-gap},
\[
 Jmt\gg X^\epsilon,
 \qquad mt\ll X^{-\epsilon}.
\]
For all sufficiently large \(X\), the nearest integer to \(m\alpha\)
is therefore zero.  Since \(\widehat w\) is Schwartz,
\[
 J\bigl|\widehat w(Jm\alpha)\bigr|
 \ll_{A,w}J(Jmt)^{-A}
 \ll_{A,w,\epsilon}JX^{-A\epsilon}.
\]
For \(k\ne0\), one has
\(|k-m\alpha|\ge |k|/2\).  Given the requested output exponent \(A\),
choose a Schwartz decay order \(B\) with
\[
 B\ge A,
 \qquad B\nu\ge A\epsilon.
\]
Then
\[
 J\sum_{k\ne0}
 \bigl|\widehat w(J(k-m\alpha))\bigr|
 \ll_{B,w}J^{1-B}
 \ll_{A,w,\epsilon}JX^{-A\epsilon}.
\]
The zero mode is bounded with the same order \(B\), and hence also by
the right-hand side of
\eqref{eq:tpc46-smooth-kernel-rapid-decay}.  The triangle inequality in
\eqref{eq:tpc46-shift-kernel-factorization} proves
\eqref{eq:tpc46-smooth-packet-rapid-decay}.
\end{proof}

\begin{remark}[Why the vectors must be orbit independent]
\label{rem:tpc46-j-independent-vectors}
If \(v_{p,r,m}\) in \eqref{eq:tpc46-Hilbert-packet} is replaced
by an arbitrary \(v_{p,r,m,j}\), then \(K_m(\alpha)\) is replaced by
the vector-valued discrete transform
\[
 \sum_jw(j/J)v_{p,r,m,j}\e(mj\alpha),
\]
which may have full shift spectrum.  The preceding theorems continue
to hold under a separately proved uniform band-limit or rapid Fourier
decay for these vector-valued profiles; they do not hold for arbitrary
orbit-output coupling.  In particular, no smoothness of an actual TPC
projective mask is being assumed implicitly here.
\end{remark}

\begin{remark}[Amplitude gap versus energy]
\label{rem:tpc46-amplitude-energy-distinction}
The vanishing in \eqref{eq:tpc46-exact-gap-vanishing} concerns the
Fourier transform of the Hilbert-valued amplitude \(F_h\).  It gives an
orthogonal low-frequency/high-alias decomposition of the total
Parseval energy \eqref{eq:tpc46-Hilbert-shift-Parseval}.  It does not
give the same spectral gap for the scalar sequence
\(h\mapsto\|F_h\|_{\Hh}^2\): multiplication in \(h\) convolves shift
frequencies, and differences of low and alias frequencies may occupy
the intermediate band.  Consequently, the amplitude gap alone is not
an atomic-normalized Bessel estimate at a prescribed shift.
\end{remark}
